一理想迈克耳孙干涉仪在空气中工作，光源为中心频率 \(\nu_0>0\) 的近似准单色非偏振光。假定各频率分量互不相干，光强谱在 \([\nu_0-\Delta\nu/2,\nu_0+\Delta\nu/2]\) 上均匀分布，区间外为零，且 \(0<\Delta\nu\ll\nu_0\)。两臂返回到同一探测点时，单独打开任一臂的积分光强相同为 \(I_0\)。定义两臂之间光程差 \(\ell\)、时间差 \(\tau=\ell/c\)。
(1) 对频率积分，严格求总探测光强 \(I(\tau)\)；(2) 求条纹可见度、第一次完全失去可见度时的光程差，以及为何之后仍可能出现旁瓣；(3) 当移动一面反射镜距离 \(x\) 时求 \(\ell\) 的变化，得到中心条纹周期与第一个可见度零点对应的镜面位移；(4) 以 \(\lambda_0=c/\nu_0\) 和窄带谱宽 \(\Delta\lambda\) 表示一阶相干长度，说明近似的条件。
